\section{Related Work}
\label{sec:rw}
\label{sec:related}

Prior work has studied test suite completeness through structural metrics and code-centric gap analysis, and has used documentation to guide test generation. However, whether a semantic gap exists between the behaviours the developer documented and the behaviours their test suites actually validate has not been studied. In the following subsections, we discuss how each of these areas provides context to our study.

\subsection{Structural Test Measures}

Code coverage and mutation kill score are the two most widely adopted proxies for test suite quality in both research and practice~\cite{papadakis2019mutation, 10.1145/2568225.2568278, 10.1145/2635868.2635929}. Coverage measures whether source lines are executed by tests, while mutation kill score assesses whether those executions are backed by meaningful assertions~\cite{10.1145/1062455.1062530}. Mutation kill score has been shown to correlate significantly with real fault detection independently of coverage, making it the best available proxy for test suite effectiveness~\cite{10.1145/2635868.2635929}.
% this is just an assumption from the literature
Together, they are often treated as complementary signals of test suite completeness. However, other studies have shown that coverage correlates only weakly with fault detection once test suite size is controlled for~\cite{10.1145/2568225.2568271}, and mutation testing incurs significant computational cost that limits its scalability in practice. 
Fundamentally, both metrics are implementation-centric, measuring properties of the code as written rather than whether the code fulfills its documented behaviours. 
% Applied to our \texttt{emptyArray} example, both metrics report full coverage, yet the singleton behaviour documented in the \texttt{@return} block goes entirely undetected.

\subsection{Complementing Structural Testing} 

Several prior projects have studied the notion of \emph{testing gap} from a code-centric perspective. They have mainly focused on situations where tests execute code but fail to meaningfully check its behaviour. Schuler \textit{et al.}'s checked-coverage technique identifies code whose effects are not validated by the assertions in the tests~\cite{schuler2013checked}. Similarly, mutation-testing research has explored pseudo-tested methods, where mutants can be introduced without causing test failures, indicating weak or ineffective tests~\cite{vera2019comprehensive}. More recently, Maton \textit{et al.} introduced \textsc{GapGrep}, which connects checked coverage and mutation testing to identify execution gaps in test suites~\cite{maton2025tests}. 

These approaches provide insight into whether the executed code is meaningfully asserted. However, they remain fundamentally code-centric and do not consider whether tests validate the behaviours developers explicitly describe in documentation. Fraser \textit{et al.} 
recognized this limitation early, proposing behavioural adequacy as a criterion grounded in whether test executions can reproduce an accurate model of system behaviour~\cite{10.1109/ICST.2012.110}. 
Despite this early recognition, to the best of our knowledge, no work has characterized how large this gap is between documented intent and what existing developer-written test suites actually validate. We address this gap empirically, using LLMs as a measurement instrument rather than a generation engine, to quantify the lower bound on how much of the documented intent of real-world Java libraries goes unvalidated in practice.

%%% i am merging this to the section Unit test section by compressing the info

% \subsection{Large Language Models} The use of Large Language Models (LLMs) in software engineering has increased significantly~\cite{hou2024largelanguagemodelssoftware}. LLMs have the ability to understand complex code and produce semantically coherent program snippets because they were trained on a vast collection of text and source code~\cite{vaswani2017attention, ouyang2022training}. 

% LLMs have been applied to tasks that bridge natural language and code. Endres \textit{et al.} explored whether LLMs can transform natural language intent (i.e., Javadoc) into formal method postconditions, showing that natural language specifications can be translated into 
% executable contracts~\cite{Endres2023CanLL}. Macke and Doyle showed that providing an LLM with incorrect documentation ``can greatly hinder code understanding'', highlighting that LLMs use documentation when reasoning about code~\cite{macke2024testingeffectcodedocumentation}. Tools like DocPrism~\cite{xu2025docprismlocalcategorizationexternal} and C4RLLaMA~\cite{11029963} use LLMs to detect inconsistencies between code and its documentation. However, existing LLM-based approaches still focus on generating new tests or oracles rather than evaluating the existing test suite. In this work, we use LLMs as an instrument to study the semantic gap between documented intent and existing 
% test suite validation across ten real-world Java projects.

\subsection{Automated Test Generation}


% Unit testing is essential for verifying software correctness. But writing comprehensive unit tests manually is slow and tedious, especially for large code bases. Because of this, many tools have been developed to automatically create tests. 
% Many automated test generation tools have been developed to improve structural coverage. These includes search based tools (e.g., EvoSuite~\cite{10.1145/2025113.2025179}), static analysis methods like symbolic execution~\cite{kurian2023automatically}, or constraint solving~\cite{blasi2022call}, dynamic analysis based methods like fuzz testers~\cite{10.1145/3611643.3616327, pacheco2007feedback}, and LLM based methods like ChatUniTest~\cite{chen2024chatunitest} or Aster~\cite{Pan2024ASTERNA}. Each of these approaches tries to generate tests that increase structural measures like code coverage.
Automated test generation tools target structural coverage through various strategies: search-based approaches (e.g., EvoSuite~\cite{10.1145/2025113.2025179}), symbolic execution~\cite{kurian2023automatically}, constraint solving~\cite{blasi2022call}, fuzzing~\cite{10.1145/3611643.3616327, pacheco2007feedback}, and LLM-based methods such as ChatUniTest~\cite{chen2024chatunitest} and ASTER~\cite{Pan2024ASTERNA}.


\begin{figure*}[ht!]
    \centering
    \includegraphics[width=0.8\textwidth,keepaspectratio]{resources/bFinderV3c.pdf}
    \caption{Process \toolname follows to identify any untested behaviours for a given documented method and the system's existing developer-written test suite.\vspace{-1em}}
    \label{fig:tool}
\end{figure*}

Some tools use method documentation to guide test generation. 
% Toradocu uses natural language parsing with lexicographic matching to translate \textit{exception}-related conditions from structured Javadoc tags into executable oracles~\cite{10.1145/2931037.2931061}. 
Toradocu extracts \textit{exception}-related conditions from structured Javadoc tags to generate executable oracles~\cite{10.1145/2931037.2931061}. 
% JDoctor extends this to extract \textit{@param}, \textit{@return}, and \textit{@throws} tags to generate precondition and postcondition contracts from structured documentation~\cite{10.1145/3213846.3213872}
JDoctor extends this approach to infer pre- and post-conditions from \textit{@param}, \textit{@return}, and \textit{@throws} tags~\cite{10.1145/3213846.3213872}, while Judot uses these inferred contracts to guide search-based test generation~\cite{denaro2025automatedtestgenerationprogram}. 
However, JDoctor's reliance on structured tags limits its reach to methods with 
formal documentation, leaving behaviours expressed in narrative text outside its scope. DISTINCT uses LLMs to generate tests by using natural language descriptions and branch-level coverage analysis~\cite{zhang2025frameworkcreatingnonregressivetest}. 

%%%% compressed LLM part from section III (large language model)%%%%
More recently, LLMs have been applied to tasks that bridge natural language and code. Endres~\textit{et al.}~\cite{Endres2023CanLL} showed that Javadoc can be translated into executable postconditions, and tools like DocPrism~\cite{xu2025docprismlocalcategorizationexternal} use LLMs to detect code-documentation inconsistencies. 
DocPrism showed that although documentation is sometimes viewed as unreliable, documentation errors are relatively infrequent ($\sim$10\%) and are typically detectable~\cite{xu2025docprismlocalcategorizationexternal}, suggesting documentation can be relied upon in practice.
%making it a reasonable source of truth for well-maintained libraries.
%%%% end compressed part

Despite this progress, these tools treat documentation as a blueprint for creating new tests or oracles from scratch, but do not evaluate whether the existing developer-written test suite already covers the documented behaviours. As a result, even with documentation-aware generation, the extent to which existing test suites fail to validate their own documented intent remains an open empirical question. 
In this paper we empirically measure behavioural gaps in both developer-written test suites and in suites created by automated test generation approaches.

